\subsection{缓增广义函数}

定义 4. --- 我们把 \(\mathbf{R}^{n}\) 上的缓增广义函数空间称为赋以有界收敛拓扑的 \(\mathcal{S}(\mathbf{R})\) 的对偶空间。记作 \(\mathcal{S}'(\mathbf{R}^{n})\) 。

由于 \(\mathcal{S}(\mathbf{R}^n)\) 是有界型的，空间 \(\mathcal{S}'(\mathbf{R}^n)\) 是完备的（TVS, III, p. 24，推论 1）。由于 \(\mathcal{S}(\mathbf{R}^n)\) 是蒙泰尔空间，\(\mathcal{S}'(\mathbf{R}^n)\) 亦然（EVT, IV, p. 19，命题 9）。于是空间 \(\mathcal{S}'(\mathbf{R}^n)\) 是自反的（EVT, IV, p. 16，定理 2）。

引理 9. --- 从 \(\mathcal{S}(\mathbf{R}^{n})\) 到 \(\mathbf{C}\) 中的线性映射 f 是缓增广义函数当且仅当：对每个族 \((\mathrm{M}_{\alpha,k})_{(\alpha,k)\in\mathbf{N}^{n}\times\mathbf{N}}\) （\(\mathbf{R}_{+}\) 中的），线性形式 f 在由这样的函数 \(\varphi\in\mathcal{S}(\mathbf{R}^{n})\) ——满足 \(\widetilde{q}_{\alpha,k}(\varphi)\leqslant\mathrm{M}_{\alpha,k}\) （对所有 \((\alpha,k)\in\mathbf{N}^{n}\times\mathbf{N}\) ）——所成之集上有界。

事实上，空间 \(\mathcal{S}(\mathbf{R}^{n})\) 是有界型的（EVT, III, p. 12，命题 2），且 \(\mathcal{S}(\mathbf{R}^{n})\) 的每个有界子集都含于陈述中所描述的有界集之一。

引理 10. --- 设 F 是 \(\mathcal{S}'(\mathbf{R}^n)\) 上具有可数基或包含一个单有界集合的滤子。那么 F 收敛于某个缓增广义函数当且仅当 \(\langle f, \varphi \rangle\) 对每个施瓦茨函数 \(\varphi\) 依 F 收敛。

特别地，缓增广义函数序列 \((f_{m})_{m\in\mathbf{N}}\) 收敛于缓增广义函数 f 当且仅当 \(\langle f_{m},\varphi\rangle\to\langle f,\varphi\rangle\) 对所有 \(\varphi\in\mathcal{S}(\mathbf{R}^{n})\) 成立。

由于 \(\mathcal{S}'(\mathbf{R}^n)\) 是弗雷歇空间，故为桶型空间（EVT, III, p. 25，命题 2 的推论），引理由巴拿赫–斯坦豪斯定理得出（EVT, III, p. 26，定理 1 的推论 3）。

典范单射 \(j\) ——\(\mathcal{D}(\mathbf{R}^{n})\) 到 \(\mathcal{S}(\mathbf{R}^{n})\) 中的——是连续的且其像稠密（IV, p. 212，引理 12），因而 \(j\) 的转置——即缓增广义函数到子空间 \(\mathcal{D}(\mathbf{R}^{n})\) 的限制映射——是从 \(\mathcal{S}'(\mathbf{R}^{n})\) 到 \(\mathcal{D}'(\mathbf{R}^{n})\) 中的单连续线性映射。我们将通过此映射把 \(\mathcal{S}'(\mathbf{R}^{n})\) 与 \(\mathcal{D}'(\mathbf{R}^{n})\) 的一个子空间等同起来。

设 \(\alpha \in \mathbf{N}^n\) 。线性映射 \(\varphi \mapsto \partial^\alpha \varphi\) ——从 \(\mathcal{S}(\mathbf{R}^n)\) 到 \(\mathcal{S}(\mathbf{R}^n)\) 中——是连续的。因此它的转置是从 \(\mathcal{S}'(\mathbf{R}^n)\) 到 \(\mathcal{S}'(\mathbf{R}^n)\) 中的连续线性映射（EVT, IV, p. 6，推论，\(b\) ））。我们用 \(\partial^\alpha\) 表示连续线性映射 \((-1)^{|\alpha|}{}^{t}\partial^\alpha\) ——从 \(\mathcal{S}'(\mathbf{R}^n)\) 到 \(\mathcal{S}'(\mathbf{R}^n)\) 中。该定义与 IV, p. 208 中关于广义函数的定义 2 相容。

设 \(h \in \mathcal{C}^{\infty}(\mathbf{R}^{n})\) 是 \(\partial^{\alpha}h\) （对每个 \(\alpha \in \mathbf{N}^{n}\) ）都具有多项式增长的函数。连续线性映射 \(\varphi \mapsto h\varphi\) （IV, p. 214，命题 14）的转置是 \(\mathcal{S}'(\mathbf{R}^{n})\) 上的连续线性映射，记作 \(f \mapsto hf\) 。
% semantic-fragments: legacy-input-seam
\relax{}
